
\documentclass{book}
%%%%%%%%%%%%%%%%%%%%%%%%%%%%%%%%%%%%%%%%%%%%%%%%%%%%%%%%%%%%%%%%%%%%%%%%%%%%%%%%%%%%%%%%%%%%%%%%%%%%%%%%%%%%%%%%%%%%%%%%%%%%%%%%%%%%%%%%%%%%%%%%%%%%%%%%%%%%%%%%%%%%%%%%%%%%%%%%%%%%%%%%%%%%%%%%%%%%%%%%%%%%%%%%%%%%%%%%%%%%%%%%%%%%%%%%%%%%%%%%%%%%%%%%%%%%
\usepackage{amssymb}
\usepackage{amsfonts}
\usepackage{amsmath}

\setcounter{MaxMatrixCols}{10}
%TCIDATA{OutputFilter=LATEX.DLL}
%TCIDATA{Version=5.50.0.2960}
%TCIDATA{<META NAME="SaveForMode" CONTENT="1">}
%TCIDATA{BibliographyScheme=Manual}
%TCIDATA{Created=Tuesday, April 15, 2025 16:35:49}
%TCIDATA{LastRevised=Thursday, January 08, 2026 16:06:09}
%TCIDATA{<META NAME="GraphicsSave" CONTENT="32">}
%TCIDATA{<META NAME="DocumentShell" CONTENT="Standard LaTeX\Standard LaTeX Book">}
%TCIDATA{Language=American English}
%TCIDATA{CSTFile=40 LaTeX Book.cst}
%TCIDATA{ComputeGeneralSettings=0,5,5,0,0,0,1}

\newtheorem{theorem}{Theorem}
\newtheorem{acknowledgement}[theorem]{Acknowledgement}
\newtheorem{algorithm}[theorem]{Algorithm}
\newtheorem{axiom}[theorem]{Axiom}
\newtheorem{case}[theorem]{Case}
\newtheorem{claim}[theorem]{Claim}
\newtheorem{conclusion}[theorem]{Conclusion}
\newtheorem{condition}[theorem]{Condition}
\newtheorem{conjecture}[theorem]{Conjecture}
\newtheorem{corollary}[theorem]{Corollary}
\newtheorem{criterion}[theorem]{Criterion}
\newtheorem{definition}[theorem]{Definition}
\newtheorem{example}[theorem]{Example}
\newtheorem{exercise}[theorem]{Exercise}
\newtheorem{lemma}[theorem]{Lemma}
\newtheorem{notation}[theorem]{Notation}
\newtheorem{problem}[theorem]{Problem}
\newtheorem{proposition}[theorem]{Proposition}
\newtheorem{remark}[theorem]{Remark}
\newtheorem{solution}[theorem]{Solution}
\newtheorem{summary}[theorem]{Summary}
\newenvironment{proof}[1][Proof]{\noindent\textbf{#1.} }{\ \rule{0.5em}{0.5em}}
\input{tcilatex}
\begin{document}

\frontmatter
\title{{\Large Fundamentals of Physics II}\\
{\Large Wave Motion, Thermal Physics, and Optics (under construction)}}
\author{R. Todd Lines \and (%
%TCIMACRO{\TeXButton{$\copyright$}{$\copyright$}}%
%BeginExpansion
$\copyright$%
%EndExpansion
copyright, R. Todd Lines, All Rights Reserved)}
\date{07/08/2025}
\maketitle
\tableofcontents

\chapter{Preface}

This set of notes is intended to be an aid to the student. It is what I
intend to present in class. There are likely errors and mistakes, so use
these notes, but don't expect perfection. If there are things that are
confusing, please talk to me or ask questions in class.

I consulted many authors and colleagues in creating these notes. Principle
among the authors was Serway and Jewett and Haliday and Resnic. Colleagues
of special note are Kevin Kelly, Brian Pyper, Steve Tercotte, and Ryan
Nielson. I especially appreciate the PH123 students of the Fall of 2006.
They were the first to test run this sets of notes, and I appreciate their
feedback (and patience).

This work is 
%TCIMACRO{\TeXButton{$\copyright$}{$\copyright$} }%
%BeginExpansion
$\copyright$
%EndExpansion
Copyright R. Todd Lines. Permission is given to BYU-Idaho for educational
purposes. All other rights are reserved.

\mainmatter

%TCIMACRO{\TeXButton{\pagenumbering{arabic}}{\pagenumbering{arabic}}}%
%BeginExpansion
\pagenumbering{arabic}%
%EndExpansion

\part{Oscillation and Waves}

\chapter{Introduction to the Course, Simple Harmonic Motion 1.15.1}

%TCIMACRO{%
%\TeXButton{\chaptermark{Simple Harmonic Motion}}{\chaptermark{Simple Harmonic Motion}}}%
%BeginExpansion
\chaptermark{Simple Harmonic Motion}%
%EndExpansion

%TCIMACRO{%
%\TeXButton{Fundamental Concepts}{\vspace{0.10in}\hspace{-0.37in}{\Large Fundamental Concepts\vspace{0.2in}}}}%
%BeginExpansion
\vspace{0.10in}\hspace{-0.37in}{\Large Fundamental Concepts\vspace{0.2in}}%
%EndExpansion

\begin{itemize}
\item Simple Harmonic Motion

\item Frequency and Period

\item Mathematical description of Simple Harmonic Motion
\end{itemize}

\section{1 Oscillation and Simple Harmonic Motion}

Last semester, in PH121 (or Dynamics) you studied how things move. We
identified a moving object (I often refer to this object as the \emph{mover
object})\emph{\ }and other objects that exerted forces on the mover (I often
refer to these other objects as the \emph{environmental objects})\emph{.}
You learned about forces and torques which get mover objects moving. You
should remember Newton's Second Law and Newton's Second Law for rotation.

You also learned and practiced a lot of math. We will continue to use the
math you learned in PH121 this semester.

But we will go beyond what we learned in PH121 to study new types of motion,
and new objects that move.

This semester, we will start with a very special type of motion. It is the
motion that results from oscillation. We call this very special type of
motion, \emph{simple harmonic motion}.

Simple harmonic motion (SHM) means a motion that repeats in the most
special, simple way. Some characteristics of this type of motion are as
follows:

1.\qquad The motion repeats in a regular way, like a grandfather clock
pendulum swings back and forth in a set amount of time. This set amount of
time is called a period.

2.\qquad The mover object moves about a center position and is symmetric
about that position. This center position is the bottom of the swing for our
grandfather clock. The idea of symmetry just means that the pendulum reaches
the same height on both sides as it swings

3.\qquad The mover object's motion can be traced out in a position vs. time
graph. For SHM, the shape of the position vs. time graph is a sinusoid.

4.\qquad The velocity of the mover object constantly changes. It is zero at
the extreme points (the points farthest from the center, or the largest
positive and the largest negative displacements). That is where it turns
around and goes back the other way. It stops there, but just for a split
second. The velocity is largest at the equilibrium position (the center
position). If we plotted a velocity vs. time graph, the velocity would also
be a sinusoid or snake-like shape.

5.\qquad The acceleration also constantly changes. And it is also a sinusoid.

That is a lot of criteria for a motion that is supposed to be simple! Let's
take an example system to see how simple harmonic motion works. I have drawn
a mass (marked $m$) attached to a spring. And lets assume that the mass is
on a frictionless surface and there is no air drag force. If we pull the
mass so the string stretches, we would get simple harmonic motion.

\FRAME{dtbpF}{2.5261in}{1.9484in}{0pt}{}{}{Figure}{\special{language
"Scientific Word";type "GRAPHIC";maintain-aspect-ratio TRUE;display
"USEDEF";valid_file "T";width 2.5261in;height 1.9484in;depth
0pt;original-width 1.2834in;original-height 0.9833in;cropleft "0";croptop
"1";cropright "1";cropbottom "0";tempfilename
'T8KIU900.wmf';tempfile-properties "XPR";}}The center position for our mass
is called the Equilibrium position, and for convenience we often define it
as the origin of our coordinate system.

\begin{center}
\rule{300pt}{1pt}

{\small Equilibrium Position: The position of the mover mass (not the
spring, which is an environmental object acting on the mover mass) when the
spring is neither stretched nor compressed.}

\rule{300pt}{1pt}
\end{center}

Let's draw a picture of simple harmonic motion. First we need a device to do
this. Suppose you go to the supermarket\footnote{%
Like Albertsons, where they still have checkers.}. But instead of putting
your ramen and granola bars on the belt at the checkout counter, you strap
on a device like this\footnote{%
Please don't really try this at the local supermarkets. They don't seem to
have any sense scientife inquarey or even a sense of humor about such things
at all.}\FRAME{dhF}{4.3474in}{1.5904in}{0pt}{}{}{Figure}{\special{language
"Scientific Word";type "GRAPHIC";maintain-aspect-ratio TRUE;display
"USEDEF";valid_file "T";width 4.3474in;height 1.5904in;depth
0pt;original-width 4.2964in;original-height 1.5549in;cropleft "0";croptop
"1";cropright "1";cropbottom "0";tempfilename
'T8KIU901.wmf';tempfile-properties "XPR";}}You can see the mover object on
the spring. But we have placed a pen on the object and that pen is tracing
out a pattern as the belt moves. You will recognize this pattern as a trig
function. The result might look something like this if you removed the belt.%
\FRAME{dhF}{3.9418in}{1.4226in}{0pt}{}{}{Figure}{\special{language
"Scientific Word";type "GRAPHIC";maintain-aspect-ratio TRUE;display
"USEDEF";valid_file "T";width 3.9418in;height 1.4226in;depth
0pt;original-width 3.8934in;original-height 1.388in;cropleft "0";croptop
"1";cropright "1";cropbottom "0";tempfilename
'T8KIU902.wmf';tempfile-properties "XPR";}}Of course, what we have made is a
position vs. time graph. We remember these from PH121! This gives us a
record of the past motion of the object. \FRAME{dtbpFX}{3.9583in}{1.8542in}{%
0pt}{}{}{Plot}{\special{language "Scientific Word";type "MAPLEPLOT";width
3.9583in;height 1.8542in;depth 0pt;display "USEDEF";plot_snapshots
TRUE;mustRecompute FALSE;lastEngine "MuPAD";xmin "0";xmax "4.7";xviewmin
"0";xviewmax "4.7";yviewmin "-10";yviewmax
"10";viewset"XY";rangeset"X";plottype 4;labeloverrides 3;x-label "t
(s)";y-label "x (cm)";axesFont "Times New
Roman,12,0000000000,useDefault,normal";numpoints 100;plotstyle
"patch";axesstyle "normal";axestips FALSE;xis \TEXUX{t};var1name
\TEXUX{$t$};function \TEXUX{$5\cos \left( 2t+0\right) $};linecolor
"blue";linestyle 1;pointstyle "point";linethickness 1;lineAttributes
"Solid";var1range "0,4.7";num-x-gridlines 100;curveColor
"[flat::RGB:0x000000ff]";curveStyle "Line";VCamFile
'T8KIU900.xvz';valid_file "T";tempfilename
'T8KIU903.wmf';tempfile-properties "XPR";}}where in this graph, $x_{\max }=5%
\unit{cm}.$ Having the power of mathematics, we know we can write an
equation that would describe this curve. From your Trigonometry experience,
we can guess that an equation for our mover object motion might look
something like%
\begin{equation*}
x\left( t\right) =x_{\max }\cos \left( \theta \right)
\end{equation*}%
Notice that the instantaneous position, $x\left( t\right) $ has to be a
function of time. Back in trigonometry we would have said a cosine function
was a function of an angle, $\theta .$ But we know this is a position vs.
time graph, so our angle must be different for different times. Let's write 
\begin{equation*}
\theta \left( t\right) =\omega t
\end{equation*}%
as strange sort of an angle. This \textquotedblleft angle\textquotedblright\
is a function of time, and let's say how fast our \textquotedblleft
angle\textquotedblright\ is changing is given by 
\begin{equation*}
\frac{d\theta }{dt}=\omega
\end{equation*}%
This is a sort of speed for how fast our angle is changing.

Using our rate of angle change, we can write our equation for simple
harmonic motion as 
\begin{equation*}
x\left( t\right) =x_{\max }\cos \left( \omega t\right)
\end{equation*}%
In the next figure $x\left( t\right) =x_{\max }\cos \left( \omega t\right) $
is plotted with two different values for $\omega .$ \FRAME{dtbpFX}{3.9583in}{%
1.3578in}{0pt}{}{}{Plot}{\special{language "Scientific Word";type
"MAPLEPLOT";width 3.9583in;height 1.3578in;depth 0pt;display
"USEDEF";plot_snapshots TRUE;mustRecompute FALSE;lastEngine "MuPAD";xmin
"0";xmax "4.7";xviewmin "0";xviewmax "4.7";yviewmin "-10";yviewmax
"10";viewset"XY";rangeset"X";plottype 4;labeloverrides 3;x-label "t
(s)";y-label "x (cm)";axesFont "Times New
Roman,12,0000000000,useDefault,normal";numpoints 100;plotstyle
"patch";axesstyle "normal";axestips FALSE;xis \TEXUX{t};var1name
\TEXUX{$t$};function \TEXUX{$5\cos \left( 2t+0\right) $};linecolor
"blue";linestyle 1;pointstyle "point";linethickness 1;lineAttributes
"Solid";var1range "0,4.7";num-x-gridlines 100;curveColor
"[flat::RGB:0x000000ff]";curveStyle "Line";function \TEXUX{$5\cos \left(
4t+0\right) $};linecolor "maroon";linestyle 2;pointstyle
"point";linethickness 1;lineAttributes "Dash";var1range
"0,4.7";num-x-gridlines 100;curveColor "[flat::RGB:0x00800000]";curveStyle
"Line";VCamFile 'T8KIU909.xvz';valid_file "T";tempfilename
'T8KIU904.wmf';tempfile-properties "XPR";}}We can see what $\omega $ does
for us. It stretches out or compresses our curve. In the blue (dashed) curve
our angle changes quickly. In the red (solid) curve the angle changes
slowly. Note that we are not plotting position vs. angle. Both plots would
be the same if we did! We are plotting position vs. time. And for the blue
curve our mass is oscillating much more quickly than it is for the red
curve. The quantity $\omega $ tells us something about how fast our mover
object is oscillating.

The quantity $x_{\max }$ is the maximum displace of our mover object from
the equilibrium position. This maximum displacement from equilibrium has a
name. It is called the \emph{amplitude}. In our equation I\ have used $%
x_{\max }$ where most books use the letter $A$ for amplitude\emph{\ }to
emphasize that the amplitude is the maximum displacement of the mover mass
from the mover mass equilibrium position. In our coordinate system, the mass
is going back and forth in the $x$ direction. Since the amplitude means the
maximum displacement, $x_{\max }$ is a good way to write amplitude. Here is
the same graph but with $x_{\max }=9\unit{cm}$\FRAME{dtbpFX}{3.9583in}{%
1.3578in}{0pt}{}{}{Plot}{\special{language "Scientific Word";type
"MAPLEPLOT";width 3.9583in;height 1.3578in;depth 0pt;display
"USEDEF";plot_snapshots TRUE;mustRecompute FALSE;lastEngine "MuPAD";xmin
"0";xmax "4.7";xviewmin "0";xviewmax "4.7";yviewmin "-10";yviewmax
"10";viewset"XY";rangeset"X";plottype 4;labeloverrides 3;x-label "t
(s)";y-label "x (cm)";axesFont "Times New
Roman,12,0000000000,useDefault,normal";numpoints 100;plotstyle
"patch";axesstyle "normal";axestips FALSE;xis \TEXUX{t};var1name
\TEXUX{$t$};function \TEXUX{$9\cos \left( 2t+0\right) $};linecolor
"blue";linestyle 1;pointstyle "point";linethickness 1;lineAttributes
"Solid";var1range "0,4.7";num-x-gridlines 100;curveColor
"[flat::RGB:0x000000ff]";curveStyle "Line";function \TEXUX{$9\cos \left(
4t+0\right) $};linecolor "maroon";linestyle 2;pointstyle
"point";linethickness 1;lineAttributes "Dash";var1range
"0,4.7";num-x-gridlines 100;curveColor "[flat::RGB:0x00800000]";curveStyle
"Line";VCamFile 'T8KIU908.xvz';valid_file "T";tempfilename
'T8KIU905.wmf';tempfile-properties "XPR";}}The two graphs for the two
different $\omega $ values are not more stretched out in time, but now they
are taller along the position axis. This means that the mass is moving
farther from the center as it oscillates.

You will remember from your trigonometry class that the \emph{period,} $T,$
tells us the time it takes for the oscillation to go through a complete
cycle. A complete cycle is when the object, say, goes from $x_{\max }$ to $%
-x_{\max }$ and then back to $x_{\max }.$ You can probably guess that how
long it takes to oscillate and how often it oscillates would be related. How
often the oscillator completes a cycle is called the frequency. The longer
the period, the lower the frequency. 
\begin{equation*}
f=\frac{1}{T}
\end{equation*}%
Think of cars passing you on your way to class. Period is like how long you
wait in between cars. Frequency is like how often cars pass. If you wait
less time between cars, the cars pass more frequently. And that is just what
our equation says! We can see from our graphs that our stretching quantity $%
\omega ,$ must be related to the frequency of our oscillation. If the
frequency is high, then $\omega $ must be large so that we reach different
\textquotedblleft angles\textquotedblright\ faster. But for the cosine
function to work we need angle units. \emph{We will choose radians} for our
units and we will write our stretching quantity as 
\begin{equation*}
\omega =2\pi f
\end{equation*}%
This works! If $\omega $ is bigger then our oscillation happens more
frequently. The $2\pi $ has units of radians. So $\omega $ has units of $%
\unit{rad}\unit{Hz}$ or more commonly $\unit{rad}/\unit{s}.$ That matches
our derivative above. Let's give a name to the quantity $\omega .$ Since it
has radians in it we might guess that it has something to do with circular
motion (more on this later) and it has frequency in it. So we will call $%
\omega $ the \emph{angular frequency}.

\section{Velocity and Acceleration}

So far we have guessed the descriptive equation for SHM. 
\begin{equation}
x\left( t\right) =x_{\max }\cos \left( \omega t\right)
\end{equation}%
This gives us the position of the particle at any given time. Since knowing
the position as a function of time is a good description of the motion of
our mover mass, we can call this equation the \emph{equation of motion} for
our mass. We will see soon that this equation is correct. But let's pretend
that we are earlier researchers and we just know the equation makes the
right shape. Still, we can learn much from knowing this equation. Since we
know how to take derivatives, and we know the derivative of position with
respect to time is the velocity, we can see that the velocity of the mass at
any given time is given by

\begin{equation}
v\left( t\right) =\frac{dx\left( t\right) }{dt}=-\omega x_{\max }\sin \left(
\omega t\right)
\end{equation}%
From our fond memories of our trigonometry class we know the maximum of a
sine function is always $1.\strut $\FRAME{dtbpFX}{4.4996in}{1.0421in}{0pt}{}{%
}{Plot}{\special{language "Scientific Word";type "MAPLEPLOT";width
4.4996in;height 1.0421in;depth 0pt;display "USEDEF";plot_snapshots
TRUE;mustRecompute FALSE;lastEngine "MuPAD";xmin "-10.001000";xmax
"10.001000";xviewmin "-10.001000";xviewmax "10.001000";yviewmin
"-1.000187";yviewmax "1.000187";viewset"XY";rangeset"X";plottype
4;labeloverrides 3;x-label "x";y-label "sin(x)";axesFont "Times New
Roman,12,0000000000,useDefault,normal";numpoints 100;plotstyle
"patch";axesstyle "normal";axestips FALSE;xis \TEXUX{x};var1name
\TEXUX{$x$};function \TEXUX{$\sin x$};linecolor "blue";linestyle
1;pointstyle "point";linethickness 1;lineAttributes "Solid";var1range
"-10.001000,10.001000";num-x-gridlines 100;curveColor
"[flat::RGB:0x000000ff]";curveStyle "Line";function \TEXUX{$1$};linecolor
"black";linestyle 2;pointstyle "point";linethickness 1;lineAttributes
"Dash";var1range "-10.001000,10.001000";num-x-gridlines 100;curveColor
"[flat::RGB:0000000000]";curveStyle "Line";function \TEXUX{$-1$};linecolor
"black";linestyle 2;pointstyle "point";linethickness 1;lineAttributes
"Dash";var1range "-10.001000,10.001000";num-x-gridlines 100;curveColor
"[flat::RGB:0000000000]";curveStyle "Line";VCamFile
'T8KIU907.xvz';valid_file "T";tempfilename
'T8KIU906.wmf';tempfile-properties "XPR";}}Notice that the $\omega $ and the 
$x_{\max }$ aren't changing for our oscillation mover mass. They are
constant. Then if we want the maximum speed we can simply set the $\sin
\left( \omega t\right) =1.$ Then the maximum speed will be 
\begin{equation}
v_{\max }=\omega x_{\max }\left( 1\right)
\end{equation}%
Notice that this does not tell us when the speed is maximum. Just what the
maximum speed is. We then have an equation for the speed of our object as a
function of time 
\begin{equation*}
v\left( t\right) =-v_{\max }\sin \left( \omega t\right)
\end{equation*}%
We will often use this trick of knowing the maximum of sine is one.

We can also find the acceleration. We just take another derivative.%
\begin{eqnarray*}
a\left( t\right) &=&\frac{dv}{dt} \\
&=&\frac{d^{2}x\left( t\right) }{dt^{2}} \\
&=&-\omega ^{2}x_{\max }\cos \left( \omega t\right)
\end{eqnarray*}%
This is the acceleration of the object attached to the spring. It's also
true that the maximum for a cosine function is $1,$ so the maximum
acceleration would be

\begin{equation}
a_{\max }=\omega ^{2}x_{\max }\left( 1\right)
\end{equation}%
so we could write our instantaneous acceleration as 
\begin{equation*}
a\left( t\right) =-a_{\max }\cos \left( \omega t\right)
\end{equation*}

These are significant results, so let's summarize. For a simple harmonic
oscillator, the instantaneous position, speed, and acceleration are given by%
\begin{eqnarray}
x\left( t\right) &=&x_{\max }\cos \left( \omega t\right) \\
v\left( t\right) &=&-\omega x_{\max }\sin \left( \omega t\right)  \notag \\
a\left( t\right) &=&-\omega ^{2}x_{\max }\cos \left( \omega t\right)  \notag
\end{eqnarray}

Let's plot $x\left( t\right) ,$ $v\left( t\right) ,$ and $a\left( t\right) $
for a specific case\FRAME{dtbpFX}{3.6357in}{2.4163in}{0pt}{}{}{Plot}{\special%
{language "Scientific Word";type "MAPLEPLOT";width 3.6357in;height
2.4163in;depth 0pt;display "USEDEF";plot_snapshots TRUE;mustRecompute
FALSE;lastEngine "MuPAD";xmin "-5";xmax "5";xviewmin "-5";xviewmax
"5";yviewmin "-20";yviewmax "20";viewset"XY";rangeset"X";plottype
4;labeloverrides 3;x-label "t";y-label "x";axesFont "Times New
Roman,12,0000000000,useDefault,normal";numpoints 100;plotstyle
"patch";axesstyle "normal";axestips FALSE;xis \TEXUX{t};var1name
\TEXUX{$t$};function \TEXUX{$5\cos \left( 2t+0\right) $};linecolor
"blue";linestyle 1;pointstyle "point";linethickness 3;lineAttributes
"Solid";var1range "-5,5";num-x-gridlines 100;curveColor
"[flat::RGB:0x000000ff]";curveStyle "Line";function \TEXUX{$-2\ast 5\sin
\left( 5t+0\right) $};linecolor "green";linestyle 2;pointstyle
"point";linethickness 3;lineAttributes "Dash";var1range
"-5,5";num-x-gridlines 100;curveColor "[flat::RGB:0x00008000]";curveStyle
"Line";function \TEXUX{$-\left( 2\right) ^{2}\left( 5\right) \cos \left(
2t+0\right) $};linecolor "red";linestyle 4;pointstyle "point";linethickness
3;lineAttributes "DotDash";var1range "-5,5";num-x-gridlines 100;curveColor
"[flat::RGB:0x00ff0000]";curveStyle "Line";VCamFile
'T8KIU906.xvz';valid_file "T";tempfilename
'T8KIU907.wmf';tempfile-properties "XPR";}}Red (solid) is the displacement,
green (dashed) is the velocity, and blue (dot-dashed) is the acceleration.
Note that each as a different maximum amplitude. Also note that they don't
rise and fall at the same time. We will describe this as being \emph{not in
phase}. We can see this again in the next graph. The red graph (top graph)
is still the displacement. The bottom graph is the velocity. And note that
the peaks are not at the same time.\FRAME{dtbpF}{2.9888in}{2.8409in}{0pt}{}{%
}{Figure}{\special{language "Scientific Word";type
"GRAPHIC";maintain-aspect-ratio TRUE;display "USEDEF";valid_file "T";width
2.9888in;height 2.8409in;depth 0pt;original-width 6.5639in;original-height
6.2396in;cropleft "0";croptop "1";cropright "1";cropbottom "0";tempfilename
'T8KIU908.wmf';tempfile-properties "XPR";}}The acceleration is $90\unit{%
%TCIMACRO{\U{b0}}%
%BeginExpansion
{{}^\circ}%
%EndExpansion
}$ \emph{out of phase} from the velocity. Let's think about why this would
be. Suppose we attach a mass to a spring and allow the mass to slide on a
frictionless surface.\FRAME{dtbpF}{2.4111in}{1.7461in}{0pt}{}{}{Figure}{%
\special{language "Scientific Word";type "GRAPHIC";maintain-aspect-ratio
TRUE;display "USEDEF";valid_file "T";width 2.4111in;height 1.7461in;depth
0pt;original-width 7.043in;original-height 5.0903in;cropleft "0";croptop
"1";cropright "1";cropbottom "0";tempfilename
'T8KIU909.wmf';tempfile-properties "XPR";}}Let's start by stretching the
spring by pulling the mass to the right and releasing. This is the situation
in the left hand part of the last figure.

\FRAME{dtbpF}{1.4114in}{1.1692in}{0pt}{}{}{Figure}{\special{language
"Scientific Word";type "GRAPHIC";maintain-aspect-ratio TRUE;display
"USEDEF";valid_file "T";width 1.4114in;height 1.1692in;depth
0pt;original-width 2.1084in;original-height 1.7417in;cropleft "0";croptop
"1";cropright "1";cropbottom "0";tempfilename
'T8KIU90A.wmf';tempfile-properties "XPR";}}The spring is pulling strongly on
the mass. We could draw a free body diagram for this situation. We will need
the spring force \FRAME{dtbpF}{1.5592in}{1.0957in}{0pt}{}{}{Figure}{\special%
{language "Scientific Word";type "GRAPHIC";maintain-aspect-ratio
TRUE;display "USEDEF";valid_file "T";width 1.5592in;height 1.0957in;depth
0pt;original-width 2.1577in;original-height 1.5091in;cropleft "0";croptop
"1";cropright "1";cropbottom "0";tempfilename
'T8KIU90B.wmf';tempfile-properties "XPR";}}Back in PH121 we said that
Hooke's Law is not something that is always true, but by \textquotedblleft
law\textquotedblright\ we mean a mathematical representation (an equation)
that comes from our mental model of how the universe works. In this case,
Hooke's law is an equation that comes from Hooke's model of how springs
work. It is a good model for most springs as long as we don't stretch them
too far. You remember Hooke's law tells us that the spring force is
proportional to how far we stretch or compress the spring from equilibrium $%
\left( \Delta x_{e}\right) $ and how stiff the spring is $\left( k\right) .$

\begin{eqnarray}
F_{s} &=&S=-k\Delta x_{e} \\
&=&-k\left( x-x_{e}\right)  \notag
\end{eqnarray}%
where $x_{e}$ is the \emph{equilibrium position of the mass. }If we assume
the equilibrium position is at $x=0$ then we can write our spring force as%
\begin{equation}
S=-kx
\end{equation}%
And if we write out Newton's second law we get 
\begin{eqnarray*}
F_{net_{x}} &=&-ma_{x}=-S_{ms} \\
F_{net_{y}} &=&ma_{y}=N_{mT}-W_{mE}
\end{eqnarray*}%
We can expect that $a_{y}$ should be zero because the mass won't lift off
the table in the $y$-direction. But in the $x$-direction 
\begin{equation*}
-a_{x}=-\frac{S_{ms}}{m}
\end{equation*}%
and knowing that 
\begin{equation*}
S_{ms}=-kx
\end{equation*}%
from Hooke's law, we can see that 
\begin{equation*}
-a_{x}=+\frac{kx}{m}
\end{equation*}%
or 
\begin{equation*}
a_{x}=-\frac{kx}{m}
\end{equation*}%
Since at our starting point $x$ is big our acceleration is big. Notice that
there is a minus sign. The minus sign tells us that $a_{x}$ must be to the
left.

But suppose the mass was to the left of $x_{e}=0$ \FRAME{dtbpF}{2.0643in}{%
0.9582in}{0pt}{}{}{Figure}{\special{language "Scientific Word";type
"GRAPHIC";maintain-aspect-ratio TRUE;display "USEDEF";valid_file "T";width
2.0643in;height 0.9582in;depth 0pt;original-width 2.0254in;original-height
0.9254in;cropleft "0";croptop "1";cropright "1";cropbottom "0";tempfilename
'T8KIU90C.wmf';tempfile-properties "XPR";}}Now $x<0$ so it is negative. That
makes 
\begin{equation*}
a_{x}=-\frac{kx}{m}
\end{equation*}%
positive. The acceleration always seems to point toward $x_{e}=0,$ the
equilibrium position for the mass. This is an important part of the
definition of simple harmonic motion, having an acceleration that always
points toward the equilibrium position. We call a force that makes the
acceleration point toward the equilibrium position \emph{restoring force}.

\begin{center}
{\small \rule{300pt}{1pt}}

{\small Restoring force:\ A force that is always directed toward the
equilibrium position}

{\small \rule{300pt}{1pt}}
\end{center}

Now lets look at our object a short time later \FRAME{dtbpF}{1.4096in}{%
1.1087in}{0pt}{}{}{Figure}{\special{language "Scientific Word";type
"GRAPHIC";maintain-aspect-ratio TRUE;display "USEDEF";valid_file "T";width
1.4096in;height 1.1087in;depth 0pt;original-width 1.375in;original-height
1.075in;cropleft "0";croptop "1";cropright "1";cropbottom "0";tempfilename
'T8KIU90D.wmf';tempfile-properties "XPR";}}We can see that in this position, 
$x=0$ so 
\begin{equation*}
a_{x}=\frac{k\left( 0\right) }{m}=0
\end{equation*}%
so at this position the acceleration is zero. That is because right at $x=0$
the spring is neither pushing nor pulling. There is no net force so no
acceleration when $x=0.$

\section{The Idea of Phase}

We said before that $x\left( t\right) ,$ $v\left( t\right) ,$ and $a\left(
t\right) $ are \textquotedblleft out of phase.\textquotedblright\ Let's look
at the idea of \textquotedblleft phase\textquotedblright\ more carefully.
Suppose you return to the grocery store and start your SHM\ device.\FRAME{dhF%
}{3.2292in}{1.1813in}{0pt}{}{}{Figure}{\special{language "Scientific
Word";type "GRAPHIC";maintain-aspect-ratio TRUE;display "USEDEF";valid_file
"T";width 3.2292in;height 1.1813in;depth 0pt;original-width
4.2964in;original-height 1.5549in;cropleft "0";croptop "1";cropright
"1";cropbottom "0";tempfilename 'T8KIU90E.wmf';tempfile-properties "XPR";}}%
But this time you work with a lab partner, and the lab partner tries to
start a stopwatch when you let go of the mass. But, due to having a slow
reaction time, your partner starts the clock too late. The resulting graph
looks like this\FRAME{dtbpFX}{3.9583in}{1.8542in}{0pt}{}{}{Plot}{\special%
{language "Scientific Word";type "MAPLEPLOT";width 3.9583in;height
1.8542in;depth 0pt;display "USEDEF";plot_snapshots TRUE;mustRecompute
FALSE;lastEngine "MuPAD";xmin "0";xmax "4.7";xviewmin "0";xviewmax
"4.7";yviewmin "-10";yviewmax "10";viewset"XY";rangeset"X";plottype
4;labeloverrides 3;x-label "t";y-label "x";axesFont "Times New
Roman,12,0000000000,useDefault,normal";numpoints 100;plotstyle
"patch";axesstyle "normal";axestips FALSE;xis \TEXUX{t};var1name
\TEXUX{$t$};function \TEXUX{$5\cos \left( 2t+\pi /3\right) $};linecolor
"blue";linestyle 1;pointstyle "point";linethickness 1;lineAttributes
"Solid";var1range "0,4.7";num-x-gridlines 100;curveColor
"[flat::RGB:0x000000ff]";curveStyle "Line";VCamFile
'T8KIU905.xvz';valid_file "T";tempfilename
'T8KIU90F.wmf';tempfile-properties "XPR";}}But we know that the only
difference between this graph and the one we had before is that the lab
partner was slow, so the graph is shifted on the time axes. We expect we can
use the same mathematical model for SHM, but we must need to change
something because the graph looks different.

We know from our algebra classes what a shift looks like. Take the
expression 
\begin{equation*}
y=x^{2}
\end{equation*}%
We can plot this to get a parabola\FRAME{dtbpFX}{2.9093in}{1.9407in}{0pt}{}{%
}{Plot}{\special{language "Scientific Word";type "MAPLEPLOT";width
2.9093in;height 1.9407in;depth 0pt;display "USEDEF";plot_snapshots
TRUE;mustRecompute FALSE;lastEngine "MuPAD";xmin "-5.001000";xmax
"10.001000";xviewmin "-5.001000";xviewmax "10.001000";yviewmin "0";yviewmax
"25.00250";viewset"XY";rangeset"X";plottype 4;axesFont "Times New
Roman,12,0000000000,useDefault,normal";numpoints 100;plotstyle
"patch";axesstyle "normal";axestips FALSE;xis \TEXUX{x};var1name
\TEXUX{$x$};function \TEXUX{$x^{2}$};linecolor "blue";linestyle 1;pointstyle
"point";linethickness 3;lineAttributes "Solid";var1range
"-5.001000,10.001000";num-x-gridlines 100;curveColor
"[flat::RGB:0x000000ff]";curveStyle "Line";VCamFile
'T8KIU904.xvz';valid_file "T";tempfilename
'T8KIU90G.wmf';tempfile-properties "XPR";}}A shifted parabola would be
expressed as 
\begin{equation*}
y=\left( x-\phi _{o}\right) ^{2}
\end{equation*}%
where $\phi _{o}$ is the amount of the shift. Suppose $\phi _{o}=2,$ then 
\begin{equation*}
y=\left( x-2\right) ^{2}
\end{equation*}%
and our shifted parabola looks like this\FRAME{dtbpFX}{2.8534in}{1.9035in}{%
0pt}{}{}{Plot}{\special{language "Scientific Word";type "MAPLEPLOT";width
2.8534in;height 1.9035in;depth 0pt;display "USEDEF";plot_snapshots
TRUE;mustRecompute FALSE;lastEngine "MuPAD";xmin "-5.001000";xmax
"10.00100";xviewmin "-5.001000";xviewmax "10.00100";yviewmin
"-0.003982";yviewmax "25.00490";viewset"XY";rangeset"X";plottype 4;axesFont
"Times New Roman,12,0000000000,useDefault,normal";numpoints 100;plotstyle
"patch";axesstyle "normal";axestips FALSE;xis \TEXUX{x};var1name
\TEXUX{$x$};function \TEXUX{$\left( x-2\right) ^{2}$};linecolor
"blue";linestyle 1;pointstyle "point";linethickness 3;lineAttributes
"Solid";var1range "-5.001000,10.00100";num-x-gridlines 100;curveColor
"[flat::RGB:0x000000ff]";curveStyle "Line";VCamFile
'T8KIU903.xvz';valid_file "T";tempfilename
'T8KIU90H.wmf';tempfile-properties "XPR";}}Recall that a shift like 
\begin{equation*}
y=\left( x-\phi _{o}\right) ^{2}
\end{equation*}%
will move the parabola to the right while 
\begin{equation*}
y=\left( x+\phi _{o}\right) ^{2}
\end{equation*}%
will move the parabola to the left.

We can use this idea to write our SHM expression for our situation with the
slow lab partner. 
\begin{equation*}
x\left( t\right) =x_{\max }\cos \left( \omega \left( t\pm \tau _{o}\right)
\right)
\end{equation*}%
In our case, the lab partner was late by $\tau _{o}=\allowbreak
4.\,\allowbreak 712\,4\unit{s}.$ This is not usually how we express the
shift, however. We usually distribute the $\omega $ so our equation looks
like%
\begin{eqnarray*}
x\left( t\right) &=&x_{\max }\cos \left( \omega t\pm \omega \tau _{o}\right)
\\
&=&x_{\max }\cos \left( \omega t\pm \omega \tau _{o}\right)
\end{eqnarray*}%
We usually use the symbol $\phi _{o}=\omega \tau _{o}$ so we will write our
SHM expression for position as a function of time as%
\begin{equation*}
x\left( t\right) =x_{\max }\cos \left( \omega t\pm \phi _{o}\right)
\end{equation*}

We could call $\phi _{o}$ the \emph{slow lab partner constant, }but that is
long and not very kind. So let's call $\phi _{o}$ by the name \emph{phase
constant}. It is also customary to call the entire expression in
parenthesis, $\left( \omega t\pm \phi _{o}\right) ,$ the phase of the cosine
function. This is especially used in the fields of Optics and
Electrodynamics.

From what we did before, we know that $x\left( t\right) ,$ $v\left( t\right)
,$ and $a\left( t\right) $ are out of phase, so there must be a phase
constant involved somehow. Let's look for it in the lectures that follow.

\chapter{2 Energy and Dynamics of SHM 1.15.2 1.15,3}

%TCIMACRO{\TeXButton{\chaptermark{SHM Energy}}{\chaptermark{SHM Energy}}}%
%BeginExpansion
\chaptermark{SHM Energy}%
%EndExpansion

Back in PH121 we started with position, velocity, and acceleration to
describe motion. We have done that again for a simple harmonic oscillator.
In PH121, after basic motion, we found that we could use the idea of a force
and Newton's second law to find the acceleration for our description of
motion. Then we changed view points and used the idea of energy to find
motion. The energy picture was easier in some ways. Let's try to develop and
energy picture for simple harmonic oscillators. Of course the goal in
physics is to describe the motion of the object, so we will say that we are
looking for the \emph{dynamics} of simple harmonic oscillators when we look
for the position, velocity, and acceleration as a function of time.

%TCIMACRO{%
%\TeXButton{Fundamental Concepts}{\vspace{0.10in}\hspace{-0.37in}{\Large Fundamental Concepts\vspace{0.2in}}}}%
%BeginExpansion
\vspace{0.10in}\hspace{-0.37in}{\Large Fundamental Concepts\vspace{0.2in}}%
%EndExpansion

\begin{itemize}
\item Initial Conditions

\item Energy and SHM

\item Equation of motion

\item Vertical oscillations
\end{itemize}

\section{Initial Conditions}

Usually, we need to know how we start our oscillator to solve a problem.
Let's see how this works.\FRAME{dtbpF}{1.5833in}{1.2187in}{0pt}{}{}{Figure}{%
\special{language "Scientific Word";type "GRAPHIC";maintain-aspect-ratio
TRUE;display "USEDEF";valid_file "T";width 1.5833in;height 1.2187in;depth
0pt;original-width 1.585in;original-height 1.2134in;cropleft "0";croptop
"1";cropright "1";cropbottom "0";tempfilename
'T8KIU90I.wmf';tempfile-properties "XPR";}}

Suppose we start the motion of a mass attached to a spring (a harmonic
oscillator) by pulling the mass to $x=x_{\max }$ and releasing it at $t=0.$
Let's see if we can find the phase. Our initial conditions require 
\begin{eqnarray*}
x\left( 0\right) &=&x_{\max } \\
v\left( 0\right) &=&0
\end{eqnarray*}

Using our equation for $x\left( t\right) $ and $v\left( t\right) $ we have%
\begin{eqnarray*}
x\left( 0\right) &=&x_{\max }=x_{\max }\cos \left( 0+\phi _{o}\right) \\
v\left( 0\right) &=&0=-v_{\max }\sin \left( 0+\phi _{o}\right)
\end{eqnarray*}%
If we choose $\phi _{o}=0,$ these conditions are met.

Notice that we needed to know the starting time and the position and the
velocity at that time. These are what we call \emph{initial conditions}. It
is still true that $x\left( t\right) $ and $v\left( t\right) $ are out of
phase. But we found $\phi _{o}=0.$ There is another phase term hiding in our
expression for $v\left( t\right) $ and to find it we need a small trig
identity.%
\begin{eqnarray*}
\sin \left( \alpha +\phi _{o}\right) &=&\cos \left( \alpha -\left( \frac{\pi 
}{2}-\phi _{o}\right) \right) \\
&=&\cos \left( \alpha -\frac{\pi }{2}+\phi _{o}\right)
\end{eqnarray*}%
so we could write 
\begin{eqnarray*}
v\left( t\right) &=&-\omega x_{\max }\sin \left( \omega t+\phi _{o}\right) \\
&=&-\omega x_{\max }\cos \left( \alpha -\frac{\pi }{2}+\phi _{o}\right)
\end{eqnarray*}%
and we can see that there was a phase shift of $-\pi /2$ hiding in our sine
function. So $v\left( t\right) $ really must be out of phase with $x\left(
t\right) .$

\subsection{A second example}

\FRAME{dtbpF}{1.7039in}{1.4422in}{0pt}{}{}{Figure}{\special{language
"Scientific Word";type "GRAPHIC";maintain-aspect-ratio TRUE;display
"USEDEF";valid_file "T";width 1.7039in;height 1.4422in;depth
0pt;original-width 1.7074in;original-height 1.4413in;cropleft "0";croptop
"1";cropright "1";cropbottom "0";tempfilename
'T8KIU90J.wmf';tempfile-properties "XPR";}}

Using the same equipment, let's start with 
\begin{eqnarray*}
x\left( 0\right) &=&0 \\
v\left( 0\right) &=&v_{i}
\end{eqnarray*}%
that is, the mover object is already moving when we start our experiment,
and we start watching just as it passes the equilibrium point.%
\begin{eqnarray*}
x\left( 0\right) &=&0=x_{\max }\cos \left( 0+\phi _{o}\right) \\
v\left( 0\right) &=&v_{i}=-v_{\max }\sin \left( 0+\phi _{o}\right)
\end{eqnarray*}

from the first equation we have 
\begin{equation*}
0=x_{\max }\cos \left( \phi _{o}\right)
\end{equation*}%
and that gives us 
\begin{equation*}
\phi _{o}=\cos ^{-1}\left( 0\right)
\end{equation*}%
so 
\begin{equation*}
\phi _{o}=\pm \frac{\pi }{2}
\end{equation*}%
really this is 
\begin{equation*}
\phi _{o}=\pm \frac{\pi }{2}\pm n\pi \qquad n=0,1,2,\cdots
\end{equation*}%
This gives the red dots in the next plot. \FRAME{dtbpFX}{4.4996in}{1.1736in}{%
0pt}{}{}{Plot}{\special{language "Scientific Word";type "MAPLEPLOT";width
4.4996in;height 1.1736in;depth 0pt;display "USEDEF";plot_snapshots
TRUE;mustRecompute FALSE;lastEngine "MuPAD";xmin "-15.001000";xmax
"15.001000";xviewmin "-15.001000";xviewmax "15.001000";yviewmin
"-1";yviewmax "1";viewset"XY";rangeset"X";plottype 4;axesFont "Times New
Roman,12,0000000000,useDefault,normal";numpoints 100;plotstyle
"patch";axesstyle "normal";axestips FALSE;xis \TEXUX{x};var1name
\TEXUX{$x$};function \TEXUX{$\cos \left( x\right) $};linecolor
"black";linestyle 1;pointstyle "point";linethickness 1;lineAttributes
"Solid";var1range "-15.001000,15.001000";num-x-gridlines 100;curveColor
"[flat::RGB:0000000000]";curveStyle "Line";function \TEXUX{$\left[
\MATRIX{2,6}{c}\VR{,,c,,,}{,,c,,,}{,,,,,}\HR{,,,,,,}\CELL{\frac{\pi
}{2}+0}\CELL{0}\CELL{\frac{\pi }{2}+1\pi }\CELL{0}\CELL{\frac{\pi }{2}+2\pi
}\CELL{0}\CELL{\frac{\pi }{2}+3\pi }\CELL{0}\CELL{\frac{\pi }{2}+4\pi
}\CELL{0}\CELL{\frac{\pi }{2}+5\pi }\CELL{0}\right] $};linecolor
"blue";linestyle 1;pointplot TRUE;pointstyle "point";linethickness
1;lineAttributes "Solid";curveColor "[flat::RGB:0x000000ff]";curveStyle
"Point";function \TEXUX{$\left[
\MATRIX{2,6}{c}\VR{,,c,,,}{,,c,,,}{,,,,,}\HR{,,,,,,}\CELL{-\frac{\pi
}{2}+0}\CELL{0}\CELL{-\frac{\pi }{2}-1\pi }\CELL{0}\CELL{\frac{\pi }{2}+2\pi
}\CELL{0}\CELL{-\frac{\pi }{2}-3\pi }\CELL{0}\CELL{-\frac{\pi }{2}-4\pi
}\CELL{0}\CELL{-\frac{\pi }{2}-5\pi }\CELL{0}\right] $};linecolor
"blue";linestyle 1;pointplot TRUE;pointstyle "point";linethickness
1;lineAttributes "Solid";curveColor "[flat::RGB:0x000000ff]";curveStyle
"Point";VCamFile 'T8KIU902.xvz';valid_file "T";tempfilename
'T8KIU90K.wmf';tempfile-properties "XPR";}}Notice that at each of these
locations $\cos \left( \phi \right) $ is zero. but let's make an agreement
that we will choose the smallest value for $\phi _{o}$ that makes $\cos
\left( \phi _{o}\right) =0.$ That is our%
\begin{equation*}
\phi _{o}=\pm \frac{\pi }{2}
\end{equation*}%
But positive and negative $\pi /2$ are the same \textquotedblleft
smallness.\textquotedblright\ We don't know the sign. Using our initial
velocity condition will let us determine the sign. Let's try it 
\begin{eqnarray*}
v_{i} &=&-v_{\max }\sin \left( \pm \frac{\pi }{2}\right) \\
v_{i} &=&\mp v_{\max } \\
v_{i} &=&\mp \omega x_{\max }
\end{eqnarray*}

From this we can see 
\begin{equation*}
x_{\max }=\pm \frac{v_{i}}{\omega }
\end{equation*}

From our initial conditions we know the initial velocity is positive, and we
insist on having positive amplitudes, so then we choose 
\begin{equation*}
\phi _{o}=-\frac{\pi }{2}
\end{equation*}%
so\ 
\begin{equation*}
\sin \left( -\frac{\pi }{2}\right) =\allowbreak -1
\end{equation*}%
and then this minus sign will cancel the one in our initial velocity
equation that to make our initial velocity positive.%
\begin{eqnarray*}
v_{i} &=&-v_{\max }\sin \left( -\frac{\pi }{2}\right) \\
&=&-v_{\max }\left( -1\right) \\
&=&v_{\max }
\end{eqnarray*}%
Our solutions are%
\begin{eqnarray*}
x\left( t\right) &=&\frac{v_{i}}{\omega }\cos \left( \omega t-\frac{\pi }{2}%
\right) \\
v\left( t\right) &=&v_{i}\sin \left( \omega t-\frac{\pi }{2}\right)
\end{eqnarray*}

Note that our solution is a set of equations! It isn't a set of numbers. We
saw this sometimes in PH121. But we will see it more in our study of
oscillations and waves. Sometimes the solutions are equations. But we did
fill in our equations with the number for the phase constant. it would be
even better if we had numbers for $\omega $ and \ $v_{i}.$

\begin{center}
{\small \rule{300pt}{1pt}}

{\small Generally to have a complete solution, you must find all the
constants based on the initial conditions. This would mean we need }$x_{\max
},${\small \ }$\omega ,${\small \ and }$\phi _{o}${\small \ to have a
complete solution.}

{\small \rule{300pt}{1pt}}
\end{center}

Let's try a complete example problem.

\subsection{Example}

A particle moving along the $x$ axis in simple harmonic motion starts from
its equilibrium position, the origin, at $t=0$ and moves to the right. The
amplitude of its motion is $2.00\unit{cm},$ and the frequency is $1.50\unit{%
Hz}.$

a) show that the position of the particle is given by%
\begin{equation*}
x=\left( 2.00\unit{cm}\right) \sin \left( 3.00\pi t\right)
\end{equation*}%
determine

b) the maximum speed and the earliest time $(t>0)$ at which the particle has
this speed,

c) the maximum acceleration and the earliest time $(t>0)$ at which the
particle has this acceleration, and

d) the total distance traveled between $t=0$ and $t=1.00\unit{s}$

\FRAME{dtbpF}{1.8775in}{0.8484in}{0pt}{}{}{Figure}{\special{language
"Scientific Word";type "GRAPHIC";maintain-aspect-ratio TRUE;display
"USEDEF";valid_file "T";width 1.8775in;height 0.8484in;depth
0pt;original-width 1.8403in;original-height 0.8164in;cropleft "0";croptop
"1";cropright "1";cropbottom "0";tempfilename
'T8KIU90L.wmf';tempfile-properties "XPR";}}

Basic Equations

\begin{eqnarray*}
x\left( t\right) &=&x_{\max }\cos \left( \omega t+\phi _{o}\right) \\
v\left( t\right) &=&-\omega x_{\max }\sin \left( \omega t+\phi _{o}\right) \\
a\left( t\right) &=&-\omega ^{2}x_{\max }\cos \left( \omega t+\phi
_{o}\right)
\end{eqnarray*}

\begin{equation*}
\omega =2\pi f
\end{equation*}

\begin{equation*}
v_{m}=\omega x_{m}
\end{equation*}

\begin{equation*}
a_{m}=\omega ^{2}x_{m}
\end{equation*}

\begin{equation*}
T=\frac{1}{f}
\end{equation*}

Variables

\begin{equation*}
\begin{tabular}{lll}
$t$ & time, initial time $=0$ & $t_{i}=0$ \\ 
$x$ & Position, Initial position =0 & $x\left( 0\right) =0$ \\ 
$v$ &  &  \\ 
$a$ &  &  \\ 
$x_{\max }$ & x amplitude & $x_{\max }=2.00\unit{cm}$ \\ 
$v_{m}$ & v amplitude &  \\ 
$a_{m}$ & a amplitude &  \\ 
$\omega $ & angular frequency &  \\ 
$\phi _{o}$ & phase constant &  \\ 
$f$ & frequency & $f=1.50\unit{Hz}$%
\end{tabular}%
\end{equation*}

Symbolic Solution

Part (a)

We can start by recognizing that we know $\omega $ because we know the
frequency. 
\begin{eqnarray*}
\omega &=&2\pi f \\
&=&2\pi \left( 1.50\unit{Hz}\right) \\
&=&9.\,\allowbreak 424\,8\unit{rad}\unit{Hz}
\end{eqnarray*}%
We also know the amplitude $x_{\max }$ which is given. 
\begin{equation*}
x_{\max }=2.00\unit{cm}
\end{equation*}%
Knowing that at $t=0$%
\begin{equation*}
x\left( 0\right) =0=x_{\max }\cos \left( 0+\phi _{o}\right)
\end{equation*}%
which we have seen before! We can guess that 
\begin{equation*}
\phi _{o}=\pm \frac{\pi }{2}
\end{equation*}

Using%
\begin{equation*}
v\left( 0\right) =-\omega x_{\max }\sin \left( 0\pm \frac{\pi }{2}\right)
\end{equation*}%
and demanding that amplitudes be positive values, and noting that at $t=0$
the velocity is positive from the initial conditions:%
\begin{equation*}
\phi =-\frac{\pi }{2}
\end{equation*}%
We also note from our trig identity that we used above%
\begin{equation*}
\cos \left( \theta -\frac{\pi }{2}\right) =\sin \left( \theta \right)
\end{equation*}%
then we have%
\begin{eqnarray*}
x\left( t\right) &=&x_{\max }\cos \left( 2\pi ft-\frac{\pi }{2}\right) \\
&=&x_{\max }\sin \left( 2\pi ft\right)
\end{eqnarray*}%
which is in the form we want. All we have to do is put in numbers. 
\begin{eqnarray*}
x\left( t\right) &=&x_{\max }\sin \left( 2\pi ft\right) \\
&=&\left( 2.00\unit{cm}\right) \sin \left( 3.00\pi t\right)
\end{eqnarray*}%
And once again our solution for a) is an equation.

Part (b)

We have a formula from our previous work for $v_{\max }$ 
\begin{eqnarray*}
v_{\max } &=&\omega x_{\max } \\
&=&2\pi fx_{\max }
\end{eqnarray*}%
so let's use it. 
\begin{eqnarray*}
v_{\max } &=&2\pi \left( 1.50\unit{Hz}\right) \left( 2.00\unit{cm}\right) \\
&=&\allowbreak 6.0\pi \unit{Hz}\unit{cm} \\
&=&\allowbreak 0.188\,50\frac{\unit{m}}{\unit{s}}
\end{eqnarray*}%
To find when $v_{\max }$ happens, take a derivative of $x\left( t\right) $%
\begin{equation*}
v\left( t\right) =v_{\max }=-\omega x_{\max }\sin \left( 2\pi ft-\frac{\pi }{%
2}\right)
\end{equation*}%
and recognize that $\sin \left( \theta \right) =1$ is at a maximum and this
happens when $\theta =\pi /2$. So we take the stuff that is in the
parenthesis for the sine function and set it equal to our angle that makes
the sine function equal to $1$.%
\begin{equation*}
\frac{\pi }{2}=2\pi ft-\frac{\pi }{2}
\end{equation*}%
We can simplify this and solve for $t$%
\begin{equation*}
\pi =2\pi ft
\end{equation*}%
\begin{equation*}
\frac{1}{2f}=t
\end{equation*}%
\begin{equation*}
t=\frac{1}{2\left( 1.50\unit{Hz}\right) }=0.333\,33\unit{s}
\end{equation*}

Part (c) Like with the velocity we must use our equations. We want
acceleration but we know it is just taking the derivative of $v\left(
t\right) $

\begin{equation*}
a\left( t\right) =-\omega ^{2}x_{\max }\cos \left( \omega t+\phi _{o}\right)
\end{equation*}%
And we recognize that the maximum is achieved when $\cos \left( \omega
t+\phi _{o}\right) =1$ or when $\omega t+\phi _{o}=0.$ We can solve this for 
$t.$

\begin{eqnarray*}
t &=&\frac{\phi _{o}}{\omega } \\
&=&\frac{-\frac{\pi }{2}}{2\pi f} \\
&=&\frac{-1}{4f} \\
&=&-0.166\,67\unit{s}
\end{eqnarray*}%
The formula for $a_{\max }$ is also already in our set of equations 
\begin{eqnarray*}
a_{\max } &=&-\omega ^{2}x_{\max } \\
&=&-(2\pi f)^{2}x_{m}
\end{eqnarray*}%
Putting in numbers gives%
\begin{eqnarray*}
a_{\max } &=&\left( 2\pi 1.5\unit{Hz}\right) ^{2}(2.00\unit{cm}) \\
&=&\allowbreak 1.\,\allowbreak 776\,5\frac{\unit{m}}{\unit{s}^{2}}
\end{eqnarray*}

Part (d)

We know the period is 
\begin{equation*}
T=\frac{1}{f}
\end{equation*}%
We should find the number of periods in $t=1.00\unit{s}$ and find the
distance traveled in one period, 
\begin{equation*}
N=\frac{t}{T}
\end{equation*}%
and multiply them together. In one period the distance traveled is 
\begin{equation*}
d=4x_{m}
\end{equation*}

\begin{equation*}
d_{tot}=d\ast N=d\ast \frac{t}{T}=4fx_{m}t
\end{equation*}%
Putting in numbers gives%
\begin{eqnarray*}
d_{tot} &=&4fx_{m}t \\
&=&4\ast 1.50\unit{Hz}\ast 2.00\unit{cm}\ast 1.00\unit{s} \\
&=&12\unit{cm} \\
&=&\allowbreak 0.12\unit{m}
\end{eqnarray*}

We have come far in only two lectures! We have a set of new equations and a
new problem type. In what we did in this lecture we used Newton's Second
Law. But in PH121 we learned that often using the idea of energy made
problems easier. Can we use energy with simple harmonic motion problems?

\section{Oscillators and Energy}

You might have noticed that we are calling mover objects that experience
simple harmonic motion (SHM) by the name \emph{simple harmonic oscillators
(SHO). }Let's consider such a SHO. Because our SHO is moving, we know there
must be energy associated with it. To understand the energy involved, let's
start with kinetic energy. Recall from PH121 that 
\begin{equation}
K=\frac{1}{2}mv^{2}
\end{equation}%
and we recall that for a spring, we have the spring potential energy given by%
\begin{equation}
U_{s}=\frac{1}{2}kx^{2}
\end{equation}%
Let's try a problem.

Using energy techniques, find the maximum velocity of the mass in terms of
the Amplitude and the angular frequency.

We want to use our problem solving steps:

\textbf{Type of problem:} This is an energy and a SHO problem:

\textbf{Drawing:} \FRAME{dtbpF}{1.8775in}{0.8484in}{0pt}{}{}{Figure}{\special%
{language "Scientific Word";type "GRAPHIC";maintain-aspect-ratio
TRUE;display "USEDEF";valid_file "T";width 1.8775in;height 0.8484in;depth
0pt;original-width 1.8403in;original-height 0.8164in;cropleft "0";croptop
"1";cropright "1";cropbottom "0";tempfilename
'T8KIU90M.wmf';tempfile-properties "XPR";}}

Variables: 
\begin{equation*}
\begin{tabular}{ll}
$E$ & energy \\ 
$K$ & Kinetic energy \\ 
$x$ & Position \\ 
$v$ & velocity or speed \\ 
$m$ & mass of the object \\ 
$k$ & spring stiffness constant%
\end{tabular}%
\end{equation*}

\textbf{Basic Equations:}

\begin{eqnarray*}
K &=&\frac{1}{2}mv^{2} \\
U_{s} &=&\frac{1}{2}kx^{2}
\end{eqnarray*}

Symbolic Answer

You might say, this is an easy problem, we know from last time that 
\begin{equation*}
v_{\max }=x_{\max }\omega
\end{equation*}%
but let's find this again using the ideas of energy. In PH121 we often found
that using energy made problems easier, so this might be worth a little more
work now. The total mechanical energy is 
\begin{equation*}
E=K+U_{s}+U_{g}
\end{equation*}%
Let's say that our oscillator is moving horizontally just like last time,
and that we define the center of mass of our object so that it's $y$%
-position is right at $y=0$ so our object has zero gravitational potential
energy. 
\begin{equation*}
E=K+U_{s}+0
\end{equation*}%
And let's say we have no friction, so that we can say that energy is
conserved. We set the initial and final energies equal to each other. 
\begin{eqnarray*}
E_{i} &=&E_{f} \\
K_{i}+U_{si} &=&K_{f}+U_{sf} \\
\frac{1}{2}mv_{i}^{2}+\frac{1}{2}kx_{i}^{2} &=&\frac{1}{2}mv_{f}^{2}+\frac{1%
}{2}kx_{f}^{2}
\end{eqnarray*}%
And let's start with our initial condition being where the spring is
stretched to $x_{\max }.$ Then at that moment, $v_{i}=0.$ And let's take our
final position when the mass is moving as fast as possible (because that is
what we are looking for, $v_{\max }$). \FRAME{dtbpF}{3.8956in}{2.6131in}{0in%
}{}{}{Figure}{\special{language "Scientific Word";type
"GRAPHIC";maintain-aspect-ratio TRUE;display "USEDEF";valid_file "T";width
3.8956in;height 2.6131in;depth 0in;original-width 3.9417in;original-height
2.6343in;cropleft "0";croptop "1";cropright "1";cropbottom "0";tempfilename
'T8KIU90N.wmf';tempfile-properties "XPR";}}We know from our PH121 experience
that this will be right when $U_{sf}=0$ (so all the energy is kinetic). Then 
\begin{equation*}
0+\frac{1}{2}kx_{\max }^{2}=\frac{1}{2}mv_{\max }^{2}+0
\end{equation*}%
or%
\begin{equation*}
\frac{1}{2}mv_{\max }^{2}=\frac{1}{2}kx_{\max }^{2}
\end{equation*}%
We can solve for $v_{\max }$ 
\begin{equation*}
v_{\max }^{2}=\frac{kx_{\max }^{2}}{m}
\end{equation*}%
\begin{equation*}
v_{\max }=x_{\max }\sqrt{\frac{k}{m}}
\end{equation*}%
But is this what we wanted? We expected that 
\begin{equation*}
v_{\max }=x_{\max }\omega
\end{equation*}%
And this is what we got so long as 
\begin{equation*}
\omega =\sqrt{\frac{k}{m}}
\end{equation*}%
And this is always true for a mass on a spring. We don't need a numeric
answer, This is reasonable (just what we expect) and the units check.

Let's look at kinetic and potential energy as a function of time. For our
Simple Harmonic Oscillator (SHO) we know the velocity as a function of time,%
\begin{equation*}
v\left( t\right) =-\omega x_{\max }\sin \left( \omega t+\phi \right)
\end{equation*}%
so the kinetic energy as a function of time must be%
\begin{eqnarray*}
K &=&\frac{1}{2}m\left( -\omega x_{\max }\sin \left( \omega t+\phi \right)
\right) ^{2} \\
&=&\frac{1}{2}m\omega ^{2}x_{\max }^{2}\sin ^{2}\left( \omega t+\phi \right)
\end{eqnarray*}%
and now we know that $\omega =\sqrt{k/m},$ so we can write the kinetic
energy as 
\begin{eqnarray*}
K &=& \\
&=&\frac{1}{2}m\left( \sqrt{\frac{k}{m}}\right) ^{2}x_{\max }^{2}\sin
^{2}\left( \omega t+\phi \right) \\
&=&\frac{1}{2}m\frac{k}{m}x_{\max }^{2}\sin ^{2}\left( \omega t+\phi \right)
\\
&=&\frac{1}{2}kx_{\max }^{2}\sin ^{2}\left( \omega t+\phi \right)
\end{eqnarray*}%
As the spring is stretched or compressed we store energy as spring potential
energy. The potential energy due to a spring acting on our SHO (mover mass)
is given by 
\begin{equation}
U_{s}=\frac{1}{2}kx^{2}
\end{equation}%
For our SHO we also know the position as a function of time 
\begin{equation*}
x\left( t\right) =x_{\max }\cos \left( \omega t+\phi _{o}\right)
\end{equation*}%
So, the potential energy as a function of time must be%
\begin{equation*}
U_{s}=\frac{1}{2}kx_{\max }^{2}\cos ^{2}\left( \omega t+\phi \right)
\end{equation*}%
\FRAME{dtbpF}{2.0906in}{0.981in}{0in}{}{}{Figure}{\special{language
"Scientific Word";type "GRAPHIC";maintain-aspect-ratio TRUE;display
"USEDEF";valid_file "T";width 2.0906in;height 0.981in;depth
0in;original-width 2.1013in;original-height 0.9704in;cropleft "0";croptop
"1";cropright "1";cropbottom "0";tempfilename
'T8KIU90O.wmf';tempfile-properties "XPR";}}Let's say, again, that our
oscillator is moving horizontally on a frictionless surface, and that we
define the center of mass of our object so that it's $y$-position is right
at $y=0$ so our object has zero gravitational potential energy%
\begin{equation*}
U_{g}=mgy=0
\end{equation*}%
so the total mechanical energy is given by%
\begin{equation*}
E=K+U_{s}
\end{equation*}%
which we can write out as%
\begin{eqnarray*}
E &=&K+U_{s} \\
&=&\frac{1}{2}kx_{\max }^{2}\sin ^{2}\left( \omega t+\phi \right) +\frac{1}{2%
}kx_{\max }^{2}\cos ^{2}\left( \omega t+\phi \right) \\
&=&\frac{1}{2}kx_{\max }^{2}\left( \sin ^{2}\left( \omega t+\phi \right)
+\cos ^{2}\left( \omega t+\phi \right) \right) \\
&=&\frac{1}{2}kx_{\max }^{2}
\end{eqnarray*}%
Where we have used the trig identity that $\sin ^{2}\left( \theta \right)
+\cos ^{2}\left( \theta \right) =1.$

It tells us that if there are no loss mechanisms (e.g. no friction) then the
energy in a harmonic oscillator never changes. And we remember that we would
say that energy is conserved for such a situation, which is not surprising
because we have already used conservation of energy for springs in PH121 and
in the previous example. But let's give a new name to a quantity that does
not change. Let's call it a \emph{constant of motion}. So we can make a
statement about the total energy for our SHO.

\begin{center}
{\small \rule{300pt}{1pt}}

{\small The total mechanical energy of an ideal SHO is a constant of motion}

{\small \rule{300pt}{1pt}}
\end{center}

If we plot the amount of kinetic and potential energy for an oscillator we
might find something like this:\FRAME{dtbpF}{4.7539in}{3.5699in}{0pt}{}{}{%
Figure}{\special{language "Scientific Word";type
"GRAPHIC";maintain-aspect-ratio TRUE;display "USEDEF";valid_file "T";width
4.7539in;height 3.5699in;depth 0pt;original-width 10.0258in;original-height
7.5247in;cropleft "0";croptop "1";cropright "1";cropbottom "0";tempfilename
'T8KIU90P.wmf';tempfile-properties "XPR";}}Note that the kinetic and
potential energy are out of phase with each other. If we plot them on the
same scale ( for the case $\phi _{o}=0$) we have\FRAME{dtbpF}{3.2655in}{%
2.2364in}{0pt}{}{}{Figure}{\special{language "Scientific Word";type
"GRAPHIC";maintain-aspect-ratio TRUE;display "USEDEF";valid_file "T";width
3.2655in;height 2.2364in;depth 0pt;original-width 3.2206in;original-height
2.1958in;cropleft "0";croptop "1";cropright "1";cropbottom "0";tempfilename
'T8KIU90Q.wmf';tempfile-properties "XPR";}}

Let's try another problem using energy of a SHO.

\section{Mathematical Representation of Simple Harmonic Motion}

We have a mathematical representation of simple harmonic motion from looking
at our graph of position vs. time. But as a problem, let's use math and what
we know from PH121 to show that our equation 
\begin{equation*}
x\left( t\right) =x_{\max }\cos \left( \omega t\right)
\end{equation*}%
must be right.

Once again recall from PH 121 
\begin{equation}
a=\frac{dv}{dt}=\frac{d^{2}x}{dt^{2}}
\end{equation}

and let's assume we have SHO moving only in the $x$-direction. \FRAME{dtbpF}{%
2.0649in}{0.9845in}{0in}{}{}{Figure}{\special{language "Scientific
Word";type "GRAPHIC";maintain-aspect-ratio TRUE;display "USEDEF";valid_file
"T";width 2.0649in;height 0.9845in;depth 0in;original-width
2.0755in;original-height 0.9739in;cropleft "0";croptop "1";cropright
"1";cropbottom "0";tempfilename 'T8KIU90R.wmf';tempfile-properties "XPR";}}%
Further assume the surface the object rests on is frictionless. Also let's
say that the equilibrium position $x_{e}=0.$ Then we can write Newton's
second law as%
\begin{eqnarray}
F_{net_{x}} &=&ma_{x}=-kx \\
m\frac{d^{2}x}{dt^{2}} &=&-kx  \notag
\end{eqnarray}%
We have a new kind of equation. If you are taking this freshman class as
a... well... freshman, you may not have seen this kind of equation before.
It is called a differential equation. The solution of this equation is a
function or functions that will describe the motion of our mass-spring
system as a function of time. It says that the way the object moves is an
equation where the second derivative is almost the same as the original
function. The only difference is some constants that are multiplied.

It is this function that we want, Let's rewrite it in terms of $\omega .$
Start by getting all the constants on the same side of the equation. 
\begin{equation*}
\frac{d^{2}x}{dt^{2}}=-\frac{k}{m}x
\end{equation*}%
We found that $\omega =\sqrt{k/m}.$then 
\begin{equation*}
\omega ^{2}=\frac{k}{m}
\end{equation*}
Then we can write our differential equation as%
\begin{equation}
\frac{d^{2}x}{dt^{2}}=-\omega ^{2}x
\end{equation}%
We need a function who's second derivative is the negative of itself with
just a constant out front. From Math 112 we know a few of these%
\begin{eqnarray*}
x\left( t\right) &=&x_{\max }\cos \left( \omega t+\phi _{o}\right) \\
x\left( t\right) &=&x_{\max }\sin \left( \omega t+\phi _{o}\right)
\end{eqnarray*}%
where $x_{\max },$ $\omega ,$ and $\phi _{o}$ are constants that we must
find. Let's choose the cosine function and explicitly take it's derivatives
to see if this function does solve our differential equation%
\begin{eqnarray}
x\left( t\right) &=&x_{\max }\cos \left( \omega t+\phi _{o}\right) \\
\frac{dx\left( t\right) }{dt} &=&-\omega x_{\max }\sin \left( \omega t+\phi
_{o}\right)  \notag \\
\frac{d^{2}x\left( t\right) }{dt^{2}} &=&-\omega ^{2}x_{\max }\cos \left(
\omega t+\phi _{o}\right)  \notag
\end{eqnarray}%
Let's substitute these expressions into our differential equation for the
motion%
\begin{eqnarray*}
\frac{d^{2}x}{dt^{2}} &=&-\omega ^{2}x \\
-\omega ^{2}x_{\max }\cos \left( \omega t+\phi _{o}\right) &=&-\omega
^{2}x_{\max }\cos \left( \omega t+\phi _{o}\right)
\end{eqnarray*}%
As long as the constant $\omega ^{2}$ is our $\omega ^{2}=k/m$ we have a
solution. We could have found that 
\begin{equation*}
\omega ^{2}=\frac{k}{m}
\end{equation*}%
by solving this differential equation, but it might not have been as
meaningful that way. We can identify $\omega $ as the angular frequency.%
\begin{equation*}
\omega =2\pi f
\end{equation*}

Thus%
\begin{equation}
\omega =\sqrt{\frac{k}{m}}=2\pi f
\end{equation}%
This says that if the spring is stiffer, we get a higher frequency, or if
the mass is larger, we get a lower frequency.

We still don't have a complete solution, because we don't know $x_{\max }$
and $\phi _{o}.$ We recognize $\phi _{o}$ as the initial phase angle. We
will have to find this by knowing the initial conditions of the motion. $%
x_{\max }$ is the amplitude. Let's look at a specific case%
\begin{equation}
\begin{tabular}{|l|}
\hline
$x_{\max }=5\unit{m}$ \\ \hline
$\phi _{o}=0$ \\ \hline
$\omega =2\unit{Hz}$ \\ \hline
\end{tabular}%
\end{equation}

\FRAME{dtbpFX}{2.7812in}{1.8542in}{0pt}{}{}{Plot}{\special{language
"Scientific Word";type "MAPLEPLOT";width 2.7812in;height 1.8542in;depth
0pt;display "USEDEF";plot_snapshots TRUE;mustRecompute FALSE;lastEngine
"MuPAD";xmin "-5";xmax "5";xviewmin "-5";xviewmax "5";yviewmin
"-10";yviewmax "10";viewset"XY";rangeset"X";plottype 4;labeloverrides
3;x-label "t";y-label "x";axesFont "Times New
Roman,12,0000000000,useDefault,normal";numpoints 100;plotstyle
"patch";axesstyle "normal";axestips FALSE;xis \TEXUX{t};var1name
\TEXUX{$t$};function \TEXUX{$5\cos \left( 2t+0\right) $};linecolor
"blue";linestyle 1;pointstyle "point";linethickness 1;lineAttributes
"Solid";var1range "-5,5";num-x-gridlines 100;curveColor
"[flat::RGB:0x000000ff]";curveStyle "Line";VCamFile
'T8KIU901.xvz';valid_file "T";tempfilename
'T8KIU90S.wmf';tempfile-properties "XPR";}}The amplitude corresponds to the
maximum displacement $x_{\max }.$ We know from trigonometry that a cosine
function has a period $T.$

\FRAME{fhF}{2.9101in}{1.9484in}{0pt}{}{}{Figure}{\special{language
"Scientific Word";type "GRAPHIC";maintain-aspect-ratio TRUE;display
"USEDEF";valid_file "T";width 2.9101in;height 1.9484in;depth
0pt;original-width 4.4616in;original-height 2.9793in;cropleft "0";croptop
"1";cropright "1";cropbottom "0";tempfilename
'T8KIU90T.wmf';tempfile-properties "XPR";}}The period is related to the
frequency%
\begin{equation}
T=\frac{1}{f}=\frac{2\pi }{\omega }
\end{equation}%
We can write the period and frequency in terms of our mass and spring
constant%
\begin{eqnarray}
T &=&2\pi \sqrt{\frac{m}{k}} \\
f &=&\frac{1}{2\pi }\sqrt{\frac{k}{m}}
\end{eqnarray}

From our graph we were able to complete our problem. But more importantly we
have two new equations for $T$ and $f$ that will help us solve more problems.

\subsection{Hanging Springs}

Let's do a problem. We now know we have forces involved with our SHOs. So
let's do a force problem. In class so far, I have always hung our masses and
springs, but we have used horizontal systems for calculations. Let's find
the equation of motion for a mass hanging from a spring.

\FRAME{dtbpF}{3.429in}{1.9977in}{0in}{}{}{Figure}{\special{language
"Scientific Word";type "GRAPHIC";maintain-aspect-ratio TRUE;display
"USEDEF";valid_file "T";width 3.429in;height 1.9977in;depth
0in;original-width 3.3831in;original-height 1.9579in;cropleft "0";croptop
"1";cropright "1";cropbottom "0";tempfilename
'T8KIU90U.wmf';tempfile-properties "XPR";}}

From observation, we would guess that we can just choose the new equilibrium
point to be $y=0$ and use the same equation%
\begin{equation*}
y\left( t\right) =y_{\max }\cos (\omega t-\phi _{o})
\end{equation*}%
let's see if that is true. Start with a free body diagram for the mass (the
hanging mass is our mover, the spring is part of the environment). There are
two forces acting on the mass. The force due to gravity, and the force due
to the spring.\FRAME{dtbpF}{0.683in}{2.1589in}{0in}{}{}{Figure}{\special%
{language "Scientific Word";type "GRAPHIC";maintain-aspect-ratio
TRUE;display "USEDEF";valid_file "T";width 0.683in;height 2.1589in;depth
0in;original-width 0.667in;original-height 2.1713in;cropleft "0";croptop
"1";cropright "1";cropbottom "0";tempfilename
'T8KIU90V.wmf';tempfile-properties "XPR";}}From our free body diagram we can
write out a Newton's second law equation%
\begin{equation*}
\Sigma F_{y}=ma_{y}=S_{ms}-W_{mE}
\end{equation*}%
We know the form of these forces%
\begin{eqnarray*}
S_{ms} &=&-k\Delta y \\
W_{mE} &=&-mg
\end{eqnarray*}%
but we need to carefully choose our origin. Let's try the top of the spring
where it attaches to the stand. If the mass just hangs there we would expect
the spring to stretch to an equilibrium length $y_{e}$ and any other motion
would either shorten or lengthen the spring. We can write our force equation
as 
\begin{eqnarray*}
ma_{y} &=&S_{ms}-W_{mE} \\
&=&-k\left( y-y_{e}\right) -mg
\end{eqnarray*}%
where $y$ can be positive or negative. If $y=0$ and the mass is just sitting
there, not oscillating, there is no acceleration. Then the stretched length
is just $y_{e}.$ In this case%
\begin{eqnarray*}
m\left( 0\right) &=&k\left( 0-y_{e}\right) -mg \\
ky_{em} &=&mg
\end{eqnarray*}%
But now let's let our mass move again. We can substitute our stationary mass
answer this into the previous equation for a moving mass%
\begin{eqnarray*}
ma &=&-k\left( y-y_{e}\right) -mg \\
&=&ky_{e}-ky-mg \\
&=&mg-ky-mg \\
&=&-ky
\end{eqnarray*}%
This gives a net force for the system of 
\begin{equation*}
F_{net}=-ky
\end{equation*}%
It is as though the system were horizontal with no gravitational force and
only a spring force. We can see that we are justified in claiming that we
could simply choose the origin at the distance $y_{em}$ from the top of the
spring, and we can use the equation 
\begin{equation*}
y\left( t\right) =y_{\max }\cos (\omega t-\phi _{o})
\end{equation*}%
as our equation of motion.

From what we have done so far, it seems like we totally redefined $\omega $
for this second physics class. But really there is a relationship between
angular velocity $\omega $ and angular frequency $\omega .$ We will explore
this in our next lecture. We will also ask what to do when there is friction.

\chapter{3 More Oscillators, Forces and Friction 1.15.4 1.15.5 1.15.6}

%TCIMACRO{\TeXButton{\chaptermark{Oscillators}}{\chaptermark{Oscillators}}}%
%BeginExpansion
\chaptermark{Oscillators}%
%EndExpansion

You have probably wondered if anyone actually uses mass-spring systems. And
we do. They were more common in the past where springs were used to store
energy ($U_{s\text{ }}$) to run clocks and toys and machines. But simple
harmonic motion can describe other systems as well. You have probably seen
an old fashioned clock with a pendulum. A pendulum almost experiences simple
harmonic motion. Let's see how this works.

%TCIMACRO{%
%\TeXButton{Fundamental Concepts}{\vspace{0.10in}\hspace{-0.37in}{\Large Fundamental Concepts\vspace{0.2in}}}}%
%BeginExpansion
\vspace{0.10in}\hspace{-0.37in}{\Large Fundamental Concepts\vspace{0.2in}}%
%EndExpansion

\begin{itemize}
\item Relationship between Simple Harmonic Motion and circular motion.

\item Pendula

\item Damping

\item Driving

\item Resonance
\end{itemize}

\section{Comparing Simple Harmonic Motion with Uniform Circular Motion}

You might have objected to our use of $\omega $ as angular frequency. Didn't 
$\omega $ mean angular speed in PH121?

That circular motion and SHM are related should not be a surprise once we
found the solutions to the equations of motion were trig functions. Recall
that the trig functions are defined on a unit circle

\FRAME{dtbpF}{2.2866in}{2.3186in}{0in}{}{}{Figure}{\special{language
"Scientific Word";type "GRAPHIC";maintain-aspect-ratio TRUE;display
"USEDEF";valid_file "T";width 2.2866in;height 2.3186in;depth
0in;original-width 2.2468in;original-height 2.2788in;cropleft "0";croptop
"1";cropright "1";cropbottom "0";tempfilename
'T8KIU90W.wmf';tempfile-properties "XPR";}}%
\begin{eqnarray}
\tan \theta &=&\frac{x}{y} \\
\cos \theta &=&\frac{x}{h} \\
\sin \theta &=&\frac{y}{h}
\end{eqnarray}

Let's relate this to our SHM equations of motion.\FRAME{dtbpF}{1.7876in}{%
1.7487in}{0pt}{}{}{Figure}{\special{language "Scientific Word";type
"GRAPHIC";maintain-aspect-ratio TRUE;display "USEDEF";valid_file "T";width
1.7876in;height 1.7487in;depth 0pt;original-width 1.7504in;original-height
1.7115in;cropleft "0";croptop "1";cropright "1";cropbottom "0";tempfilename
'T8KIU90X.wmf';tempfile-properties "XPR";}}Look at the projection $x$ of the
point $P$ on the $x$ axis. This is just the $x$-component of the position!
Lets follow this projection as $P$ travels around the circle. We find the
projection ranges from $-x_{\max }$ to $x_{\max }.$ If we watch closely we
find the projection's velocity is zero at the extreme points and is a
maximum in the middle. This projection is given as the cosine of the vector
from the origin to $P.$ It is just taking the $x$-component! This
projection, indeed fits our SHO solution.

Now lets define a projection of $P$ onto the $y$ axis. Again we have SHM,
but this time the projection is a sine function. Because 
\begin{equation}
\cos \left( \theta -\frac{\pi }{2}\right) =\sin \left( \theta \right)
\end{equation}%
we can see that this is just a SHO that is $90\unit{%
%TCIMACRO{\U{b0}}%
%BeginExpansion
{{}^\circ}%
%EndExpansion
}=\frac{\pi }{2}\unit{rad}$ out of phase. It is probably worth recalling
that a projection of one vector on another can be expressed by a dot
product. We could express our length $x$ as 
\begin{equation*}
x=\overrightarrow{\mathbf{r}}\cdot \mathbf{\hat{\imath}}=r\cos \theta
\end{equation*}%
where $r$ is the radius of the circle. \FRAME{dtbpF}{2.3687in}{2.3186in}{0in%
}{}{}{Figure}{\special{language "Scientific Word";type
"GRAPHIC";maintain-aspect-ratio TRUE;display "USEDEF";valid_file "T";width
2.3687in;height 2.3186in;depth 0in;original-width 2.3281in;original-height
2.2788in;cropleft "0";croptop "1";cropright "1";cropbottom "0";tempfilename
'T8KIU90Y.wmf';tempfile-properties "XPR";}}We can see that when $\theta =0$
we have $x=r$ and this will be the largest $x$ value, so $r=x_{\max }.$

So by projecting circular motion onto the $x$-axis 
\begin{equation*}
x\left( t\right) =x_{\max }\cos \left( \theta \right)
\end{equation*}%
But $\theta $ changes in time. We can recall from our PH121 or Dynamics
experience that the angular speed 
\begin{equation*}
\omega =\frac{\Delta \theta }{\Delta t}
\end{equation*}%
or, if we agree to start from $\theta _{i}=0$ and $t_{i}=0,$ 
\begin{equation*}
\omega =\frac{\theta }{t}
\end{equation*}%
so 
\begin{equation*}
\theta =\omega t
\end{equation*}%
then we have 
\begin{equation*}
x\left( t\right) =x_{\max }\cos \left( \omega t\right)
\end{equation*}%
Which is just our equation for SHM. Now we can see why we used
\textquotedblleft $\omega $\textquotedblright\ for both angular speed and
angular frequency. Really they are very related. Both tell us something
about how fast a cyclic event happens. For motion around a circle, one is
how fast the point $P$ goes around the circle, and the other is how often
the projection goes back and forth. It makes sense that these have to be the
same rate.\FRAME{dtbpFU}{2.5694in}{2.5201in}{0pt}{\Qcb{The projection of
circular motion onto the $x$-axis gives simple harmonic motion.}}{}{Figure}{%
\special{language "Scientific Word";type "GRAPHIC";maintain-aspect-ratio
TRUE;display "USEDEF";valid_file "T";width 2.5694in;height 2.5201in;depth
0pt;original-width 2.5287in;original-height 2.4785in;cropleft "0";croptop
"1";cropright "1";cropbottom "0";tempfilename
'T8KIU90Z.wmf';tempfile-properties "XPR";}}

Let's go back to our projection on the $y$-axis.\FRAME{dtbpF}{2.2753in}{%
2.2269in}{0pt}{}{}{Figure}{\special{language "Scientific Word";type
"GRAPHIC";maintain-aspect-ratio TRUE;display "USEDEF";valid_file "T";width
2.2753in;height 2.2269in;depth 0pt;original-width 2.2355in;original-height
2.1862in;cropleft "0";croptop "1";cropright "1";cropbottom "0";tempfilename
'T8KIU910.wmf';tempfile-properties "XPR";}}We found that we can describe
this projection as 
\begin{equation*}
y\left( t\right) =y_{\max }\sin \left( \omega t\right)
\end{equation*}

We will choose the cosine function, but from our trig experience it should
be clear that these projections are equivalent, just $90\unit{%
%TCIMACRO{\U{b0}}%
%BeginExpansion
{{}^\circ}%
%EndExpansion
}$ out of phase

\begin{equation*}
\cos \left( \theta \pm \frac{\pi }{2}\right) =\mp \sin \theta
\end{equation*}

So%
\begin{equation*}
x\left( t\right) =x_{\max }\cos \left( \omega t+\phi _{o}\right)
\end{equation*}%
could be a sine function if $\phi _{o}=\pm \pi /2.$ In this way we have
incorporated both possibilities into one function.

Note that this is just the function we guessed from our observation!

\begin{center}
{\small \rule{300pt}{1pt}}

{\small We see that uniform circular motion can be thought of as the
combination of two SHOs, with a phase difference of }$\pi /2\unit{rad}.$

{\small The angular velocity is given by}%
\begin{equation}
\omega =\frac{v}{r}
\end{equation}

{\small \rule{300pt}{1pt}}
\end{center}

\FRAME{dtbpF}{2.5997in}{2.6184in}{0pt}{}{}{Figure}{\special{language
"Scientific Word";type "GRAPHIC";maintain-aspect-ratio TRUE;display
"USEDEF";valid_file "T";width 2.5997in;height 2.6184in;depth
0pt;original-width 2.6201in;original-height 2.6388in;cropleft "0";croptop
"1";cropright "1";cropbottom "0";tempfilename
'T8KIU911.wmf';tempfile-properties "XPR";}}A particle traveling on the $x$%
-axis in SHM will travel from $x_{\max }\ $to $-x_{\max }$ and from $%
-x_{\max }\ $to $x_{\max }$ (one complete period, $T$) while the particle
traveling with $P$ makes one complete revolution. Thus, the angular
frequency $\omega $ of the SHO and the angular velocity of the particle at $%
P $ are the same. The magnitude of the tangential velocity is then%
\begin{equation}
v_{t}=\omega r=\omega x_{\max }
\end{equation}%
and the projection of this velocity onto the $x$-axis is 
\begin{equation}
v_{tx}=-\omega x_{\max }\sin \left( \omega t+\phi _{o}\right)
\end{equation}%
Which is just what we expected from our earlier observation!\FRAME{dtbpF}{%
2.4794in}{2.5079in}{0pt}{}{}{Figure}{\special{language "Scientific
Word";type "GRAPHIC";maintain-aspect-ratio TRUE;display "USEDEF";valid_file
"T";width 2.4794in;height 2.5079in;depth 0pt;original-width
2.4388in;original-height 2.4664in;cropleft "0";croptop "1";cropright
"1";cropbottom "0";tempfilename 'T8KIU912.wmf';tempfile-properties "XPR";}}

The centripetal acceleration of a particle at $P$ is given by 
\begin{equation}
a_{c}=\frac{v_{t}^{2}}{r}=\frac{v_{t}^{2}}{x_{\max }}=\frac{\omega
^{2}x_{\max }^{2}}{x_{\max }}=\omega ^{2}x_{\max }
\end{equation}%
The direction of the acceleration is inward toward the origin. Of course, we
just want the $x$-component of this, so again we make a projection. If we
project this onto the $x$-axis we have%
\begin{equation}
a_{x}=-\omega ^{2}x_{\max }\cos \left( \omega t+\phi \right)
\end{equation}%
Again this is just want we expected from our observation.

So now we have shown that our set of equations 
\begin{eqnarray}
x\left( t\right) &=&x_{\max }\cos \left( \omega t\right) \\
v\left( t\right) &=&-\omega x_{\max }\sin \left( \omega t\right)  \notag \\
a\left( t\right) &=&-\omega ^{2}x_{\max }\cos \left( \omega t\right)  \notag
\end{eqnarray}%
is correct for our harmonic oscillator.

\section{Our first problem type: Simple Harmonic Motion}

You have probably thought by now that we have a new problem type. We can
call it the \emph{simple harmonic motion }problem type or SHM. The equations
we have so far for this problem type are 
\begin{eqnarray}
x\left( t\right) &=&x_{\max }\cos \left( \omega t+\phi _{o}\right) \\
v\left( t\right) &=&-\omega x_{\max }\sin \left( \omega t+\phi _{o}\right) 
\notag \\
a\left( t\right) &=&-\omega ^{2}x_{\max }\cos \left( \omega t+\phi
_{o}\right)  \notag
\end{eqnarray}%
\begin{equation*}
\omega =2\pi f
\end{equation*}

\begin{equation*}
v_{m}=\omega x_{m}
\end{equation*}

\begin{equation*}
a_{m}=\omega ^{2}x_{m}
\end{equation*}%
\begin{equation*}
T=\frac{1}{f}
\end{equation*}

\begin{equation*}
U_{s}=\frac{1}{2}kx_{\max }^{2}\cos ^{2}\left( \omega t+\phi \right)
\end{equation*}%
\begin{equation*}
K=\frac{1}{2}m\omega ^{2}x_{\max }^{2}\sin ^{2}\left( \omega t+\phi \right)
\end{equation*}

\begin{eqnarray}
T &=&2\pi \sqrt{\frac{m}{k}} \\
f &=&\frac{1}{2\pi }\sqrt{\frac{k}{m}}
\end{eqnarray}

We have used energy to describe simple harmonic motion, and we have found
our equation of motion using a differential equation for mass-spring
systems. And we have simple harmonic motion in the $y$-direction including
the weight force due to gravity for mass-spring systems. But there are other
systems that exhibit simple harmonic motion and many others that nearly
exhibit simple harmonic motion. Let's study a few of these systems.

\section{The Simple Pendulum}

\FRAME{dtbpF}{1.7305in}{1.6737in}{0pt}{}{}{Figure}{\special{language
"Scientific Word";type "GRAPHIC";maintain-aspect-ratio TRUE;display
"USEDEF";valid_file "T";width 1.7305in;height 1.6737in;depth
0pt;original-width 1.7349in;original-height 1.6773in;cropleft "0";croptop
"1";cropright "1";cropbottom "0";tempfilename
'T8KIU913.wmf';tempfile-properties "XPR";}}A simple pendulum is a mass on a
string. The mass is called a \textquotedblleft bob.\textquotedblright\
Usually we study the motion of the pendulum bob, so let's consider the
pendulum the mover object. A simple pendulum bob exhibits periodic motion,
but not exactly simple harmonic motion. Let's set up our pendulum and
measure an angle $\theta $ starting with $\theta =0$ when the pendulum is
just hanging. Then as the pendulum moves up $\theta $ will increase.

The forces on the bob are gravity $\overrightarrow{\mathbf{W}}$, and $%
\overrightarrow{\mathbf{T}}$ the tension on the string. The tangential
component of $W$ is always directed toward $\theta =0$, the just hanging
equilibrium position. This is a restoring force! Let's call the path the bob
takes \textquotedblleft $s$\textquotedblright\ because it is an arclength.\
Then from Jr. High geometry we recall\footnote{%
Really I did not recall this, I\ have to look it up every time, but $s$ is
called the \emph{arclength}.}%
\begin{equation}
s=L\theta
\end{equation}%
We will use a the cylindrical or $rtz$ coordinate system. The radial axis is
directed along the string. The tangential direction is along the circular
path the bob takes and is always tangent to the path. In this coordinate
system, we can solve for the part of the force directed along the path. This
is the restoring part of the net force. Remember from Newton's second law
the radial and tangential components of the force are 
\begin{eqnarray*}
F_{r} &=&ma_{r} \\
F_{t} &=&ma_{t}
\end{eqnarray*}%
and 
\begin{equation*}
a_{t}=\frac{d^{2}s}{dt^{2}}
\end{equation*}%
Let's write out Newton's second law equations with the sum of the forces
part.%
\begin{eqnarray*}
-ma_{r} &=&-T+W\sin \left( 90\unit{%
%TCIMACRO{\U{b0}}%
%BeginExpansion
{{}^\circ}%
%EndExpansion
}-\theta \right) \\
ma_{t} &=&-W\cos \left( 90\unit{%
%TCIMACRO{\U{b0}}%
%BeginExpansion
{{}^\circ}%
%EndExpansion
}-\theta \right)
\end{eqnarray*}%
The second, tangential equation is the important one. It tells us how the
bob will move. Let's use another trig identity 
\begin{equation*}
\cos \left( \frac{\pi }{2}-\theta \right) =\sin \left( \theta \right)
\end{equation*}%
Then our tangential Newton's second law equation becomes%
\begin{equation*}
ma_{t}=-W\sin \left( \theta \right)
\end{equation*}%
and we can solve for the tangential acceleration, $a_{t}$ 
\begin{eqnarray*}
a_{t} &=&-\frac{W}{m}\sin \left( \theta \right) \\
&=&-g\sin \theta
\end{eqnarray*}

We have two expressions for $a_{t.}$ We can set them equal 
\begin{equation}
\frac{d^{2}s}{dt^{2}}=-g\sin \theta  \label{pendulum1}
\end{equation}%
Remember that $s=L\theta $ so we could write the left hand side as%
\begin{equation*}
\frac{d^{2}s}{dt^{2}}=\frac{d^{2}}{dt^{2}}\left( L\theta \right) =L\frac{%
d^{2}}{dt^{2}}\left( \theta \right)
\end{equation*}%
then equation (\ref{pendulum1}) becomes%
\begin{equation*}
L\frac{d^{2}}{dt^{2}}\left( \theta \right) =-g\sin \theta
\end{equation*}%
or 
\begin{equation*}
\frac{d^{2}\theta }{dt^{2}}=-\frac{g}{L}\sin \theta
\end{equation*}%
This is a differential equation much like our differential equation for a
harmonic oscillator, 
\begin{equation*}
\frac{d^{2}x}{dt^{2}}=-\omega ^{2}x
\end{equation*}%
except it has a sine function in it. But, if we take $\theta $ as a very
small angle, then 
\begin{equation}
\sin \left( \theta \right) \approx \theta  \label{small angle approximation}
\end{equation}%
We can plot this to see that it works. The red curve is $y=\sin \left(
\theta \right) $ and the blue curve is $y=\theta $ and both are plotted as a
function of $\theta .$ Notice that for less than $\theta <1\unit{rad}$ the
blue and red curves look a lot a like.

\end{document}
